%Plantilla elaborada por Dairo Andres Sierra Zamora
%Reformada por Luis García

%------------------ PREÁMBULO ------------------
\documentclass[dvipsnames, 11pt]{beamer}
\usepackage[orientation=landscape,size=custom,width=16,height=9,scale=0.5,debug]{beamerposter}
\usepackage[utf8]{inputenc}				% Acentos en español ñ y ó
\usepackage[spanish]{babel}				% Chapter -> Capítulo etc...
\usepackage[T1]{fontenc}				% Acentos como ö
\usepackage{lmodern}					% Latin modern fonts
\usepackage{amsmath,amssymb,amsfonts,latexsym,stmaryrd}
\usepackage{graphicx, fancyhdr, wrapfig, tikz, textpos}
\usetheme{Madrid}
\usecolortheme[RGB={0, 128, 128}]{structure}

\definecolor{miGrisOsc}{RGB}{75, 75, 75}
\definecolor{miGrisClaro}{RGB}{156, 156, 156}
\definecolor{miRojo1}{RGB}{255, 230, 230}
\definecolor{miRojo2}{RGB}{255, 204, 204}
\definecolor{miVerde1}{RGB}{230, 242, 230}
\definecolor{miVerde2}{RGB}{204, 229, 204}
\definecolor{miTeal}{RGB}{0, 128, 128}
\definecolor{miAmarilloFluor}{RGB}{255, 255, 0}


\usefonttheme{professionalfonts}		% Ecuaciones con letra profesional	
\setbeamersize{text margin left=10mm, text margin right=10mm}
\setbeamercolor{palette primary}{use=structure,fg=white,bg=structure.fg}
\setbeamercolor{title}{fg=white,bg=miTeal}
\setbeamercolor{frametitle}{fg=white,bg=teal}
\setbeamercolor{date in head/foot}{fg=white}
\setbeamercolor{author in head/foot}{fg=white}
\setbeamercolor{title in head/foot}{fg=white}
%\setbeamercolor{navigation symbols}{fg=teal,bg=teal}
\setbeamercolor{block title alerted}{fg=miAmarilloFluor,bg=magenta}
\setbeamercolor{block body alerted}{bg=miRojo1}

\setbeamercolor{block title example}{fg=teal,bg=miVerde2}
\setbeamercolor{block body example}{bg=miVerde1}
\setbeamercolor{bibliography item}{fg=teal}
\setbeamercolor*{bibliography entry title}{fg=miTeal}
\setbeamercolor*{bibliography entry author}{fg=magenta}
\setbeamercolor{bibliography entry location}{fg=black}
\setbeamercolor{bibliography entry note}{fg=magenta}

\setbeamertemplate{bibliography item}[text]

\setbeamertemplate{bibliography item}[text]


%------------------ INFORMACIÓN ------------------
\title[Abreviación MPU]{Modelo Presentaciones Universidad}
\subtitle[E. P]{Ejemplo Presentación}
\author[Fulanito de Tal]{Fulanito de tal \vspace{2mm}}
\institute[FUKL]{\includegraphics[scale=0.1]{figures/Logo Konrad.pdf}\vspace{5mm}\\{\large Fundación Universitaria Konrad Lorenz}\\\vspace{2mm}{\normalsize Facultad de Matemáticas e Ingeniería}}
\date{\today}
\subject{Geometría Diferencial}
% Las abreviaciones se usan para las barras de navegación y los pies de página


%------------------ OTRAS OPCIONES DEL DOCUMENTO ------------------
%\setbeamertemplate{navigation symbols}			% Oculta o muestra los símbolos de navegación
\setbeamercovered{transparent}			% Hace que lo oculto sea ligeramente viseible
%\usepackage{multicol}			% Paquetes adicionales de estilo {columnas, }
\newtheorem{teo}{Teorema}
\newtheorem{defi}{Definición}
\newtheorem{nota}{Nota}

%------------------ INICIO DEL DOCUMENTO ------------------
\begin{document}

\setbeamertemplate{background}
{
	\parbox[c][\paperheight]{\paperwidth}
	{
		\vfill 
		\begin{tikzpicture}
			\node[opacity=.1]
			{
				\includegraphics[scale=0.85]{figures/Logo Konrad Simple.pdf}
			};
		\end{tikzpicture}
		\vspace{0.5cm} \hspace{-3cm}
	}
}


	
%------------------ TITLE PAGE ------------------
\begin{frame}[plain]
	\titlepage
\end{frame}


\addtobeamertemplate{frametitle}{}{%
\begin{textblock*}{100mm}(0.86\textwidth,-1.1cm)
	\includegraphics[scale=0.08]{figures/LOGO-BASE-KONRAD-GRISES.png}
\end{textblock*}}



%------------------ NEW FRAME ------------------
\begin{frame}{Inicio}{Subtítulo}
	\begin{itemize}
		\item primero
		\item segundo 
		\item tercero
	\end{itemize}	
\end{frame}



%------------------ NEW FRAME ------------------
\begin{frame}{Título}{Subtítulo}
	\begin{block}{Importante}
		Así podemos resaltar información relevante...
	\end{block}
	
	\begin{alertblock}{Alertblock}
		contenidos...
	\end{alertblock}
	
	\begin{exampleblock}{\tt Exampleblock}
		contenidos...
	\end{exampleblock}
\end{frame}



%------------------ NEW FRAME ------------------
\begin{frame}{Multicolumnas}{Ejemplo}
	
	\begin{columns}[t]
		\begin{column}{.48\linewidth}
			contenidos...
		\end{column}
		
		\begin{column}{.48\linewidth}
			contenidos...
		\end{column}
	\end{columns}
	
\end{frame}



%------------------ NEW FRAME ------------------
\begin{frame}{Entornos theorem, proof etc}{Ejemplos}
	
	\begin{teo}[Teorema de Pitágoras]
		\begin{equation*}
			a^2 + b^2 = c^2	
		\end{equation*}
	\end{teo}
	
	\begin{proof}
		contenidos...
	\end{proof}
	
\end{frame}



%------------------ NEW FRAME ------------------
\begin{frame}{Teorema y definiciones}{Ejemplos}
	
	\begin{teo}[suma]
		$$2+2=4$$
	\end{teo}
	
	\begin{defi}[suma]
		la suma es muy importante
	\end{defi}
	
\end{frame}



%------------------ REFERENCIAS ------------------
\begin{frame}[allowframebreaks]
	\frametitle{Bibliografía}

	Algunos libros que recomiendo: \cite{Alfonso2010, Beatriz2011}

	\bibliographystyle{unsrt}
	\bibliography{biblio}
	
\end{frame}


\end{document}
%------------------ FIN DEL DOCUMENTO ------------------


